% =============================================================================
% ESL FRAMEWORK - APPENDIX F: PAPER TYPE CLASSIFICATION (PTC)
% =============================================================================
% Authors: Gerhard Fehr & Andrea Fehr
% FehrAdvice & Partners AG, Zürich
% Version: Extended Monograph
% =============================================================================

\documentclass[11pt,a4paper]{article}

% --- Packages ---
\usepackage[utf8]{inputenc}
\usepackage[T1]{fontenc}
\usepackage[english]{babel}
\usepackage{amsmath,amssymb,amsthm}
\usepackage{booktabs}
\usepackage{longtable}
\usepackage{hyperref}
\usepackage{geometry}
\usepackage{xcolor}
\usepackage{tcolorbox}
\usepackage{array}
\usepackage{multirow}
\usepackage{enumitem}
\usepackage{fancyhdr}

\geometry{margin=2.5cm}

% --- Colors ---
\definecolor{eslblue}{RGB}{0,82,147}
\definecolor{eslgreen}{RGB}{0,128,0}
\definecolor{eslorange}{RGB}{230,159,0}
\definecolor{eslred}{RGB}{204,51,17}
\definecolor{eslgray}{RGB}{128,128,128}
\definecolor{eslpurple}{RGB}{148,0,211}

% --- Boxes ---
\newtcolorbox{eslbox}[1][]{
  colback=eslblue!5!white,
  colframe=eslblue,
  title={ESL Assessment},
  fonttitle=\bfseries,
  #1
}

\newtcolorbox{keyinsight}[1][]{
  colback=eslgreen!5!white,
  colframe=eslgreen!80!black,
  title={Key Insight},
  fonttitle=\bfseries,
  #1
}

\newtcolorbox{warningbox}[1][]{
  colback=eslorange!5!white,
  colframe=eslorange!80!black,
  title={Methodological Note},
  fonttitle=\bfseries,
  #1
}

\newtcolorbox{definitionbox}[1][]{
  colback=white,
  colframe=eslblue!80!black,
  title={Definition},
  fonttitle=\bfseries,
  #1
}

\newtcolorbox{examplebox}[1][]{
  colback=eslgray!10!white,
  colframe=eslgray!80!black,
  title={Example},
  fonttitle=\bfseries,
  #1
}

% --- Theorems ---
\newtheorem{definition}{Definition}[section]
\newtheorem{principle}{Principle}[section]

% --- Header/Footer ---
\pagestyle{fancy}
\fancyhf{}
\fancyhead[L]{\small ESL Framework -- Extended Monograph}
\fancyhead[R]{\small Appendix F: Paper Type Classification}
\fancyfoot[C]{\thepage}

% =============================================================================
% TITLE
% =============================================================================

\title{%
    \Large \textbf{Empirical Stability of Language (ESL)} \\[0.5em]
    \large A Universal Quality Assurance Framework for \\
    Scientific Claim Validation \\[1em]
    \LARGE \textbf{Appendix F: Paper Type Classification (PTC)} \\[0.3em]
    \large Type-Specific Evidence Scales and Workflow Integration
}

\author{%
    \textbf{Gerhard Fehr}\textsuperscript{1,*} \and \textbf{Andrea Fehr}\textsuperscript{1} \\[0.5em]
    \textsuperscript{1}FehrAdvice \& Partners AG, Klausstrasse 4, 8008 Zürich, Switzerland \\[0.3em]
    \textsuperscript{*}Corresponding author: gerhard.fehr@fehradvice.com
}

\date{December 26, 2025}

\begin{document}

\maketitle

% =============================================================================
% ABSTRACT
% =============================================================================

\begin{abstract}
\noindent
This appendix extends the ESL framework to accommodate different types of scientific publications. While the core K-metric ($K = \min(B,E)/\max(B,E)$) remains universal, the operationalization of Evidence ($E$) varies by paper type. We introduce Paper Type Classification (PTC) as a mandatory first step in ESL assessment, define nine paper types with their characteristic claim structures, and provide type-specific evidence scales. The appendix establishes which paper types are amenable to replication probability estimation (R-Score) versus those requiring alternative quality metrics. Integration points with the main ESL workflow are specified.

\vspace{0.5em}
\noindent\textbf{Keywords:} Paper type classification, evidence scales, replication probability, meta-analysis, theoretical papers, methodology papers
\end{abstract}

\tableofcontents
\newpage

% =============================================================================
% F.1 MOTIVATION
% =============================================================================

\section{Motivation: Why Paper Type Matters}
\label{sec:motivation}

The ESL framework was initially developed for empirical studies, where ``evidence'' refers to study design quality and ``replication probability'' has a clear operational meaning. However, scientific literature encompasses diverse publication types, each with distinct claim structures and evidence standards.

\begin{keyinsight}[title={The Generalization Challenge}]
The K-metric measures \textbf{claim-evidence alignment}:
\begin{equation}
K = \frac{\min(B, E)}{\max(B, E)}
\end{equation}

This is universally applicable, but:
\begin{itemize}
    \item What counts as ``evidence'' varies by paper type
    \item Not all paper types have ``replication probability'' in the traditional sense
    \item E-scale calibration must be type-specific
\end{itemize}
\end{keyinsight}

\subsection{The Core Insight}

ESL's K-metric measures whether \textit{claim strength matches evidence strength}. This applies to:

\begin{itemize}
    \item Empirical claims (``Our study shows X'') $\rightarrow$ Evidence from study design
    \item Theoretical claims (``X logically implies Y'') $\rightarrow$ Evidence from proof rigor
    \item Synthetic claims (``The literature shows X'') $\rightarrow$ Evidence from source quality
    \item Methodological claims (``Method M works for problem P'') $\rightarrow$ Evidence from validation
\end{itemize}

Paper Type Classification (PTC) enables appropriate E-scale selection for each case.

% =============================================================================
% F.2 PAPER TYPE TAXONOMY
% =============================================================================

\section{Paper Type Taxonomy}
\label{sec:taxonomy}

We distinguish nine primary paper types based on their epistemic function and claim structure.

\begin{table}[h]
\centering
\caption{Scientific Paper Type Taxonomy}
\label{tab:taxonomy}
\begin{tabular}{clp{6cm}}
\toprule
\textbf{Code} & \textbf{Paper Type} & \textbf{Description} \\
\midrule
\textbf{EMP} & Empirical Study & New data collection and analysis \\
\textbf{META} & Meta-Analysis & Quantitative synthesis of multiple studies \\
\textbf{REV} & Review / Survey & Qualitative synthesis of existing research \\
\textbf{THE} & Theoretical Paper & New model, framework, or formal theory \\
\textbf{METH} & Methods Paper & New methodology, tool, or procedure \\
\textbf{REP} & Replication Study & Reproduction of existing study \\
\textbf{SIM} & Simulation Study & Computational modeling and analysis \\
\textbf{CASE} & Case Study & In-depth analysis of single case(s) \\
\textbf{COM} & Commentary / Letter & Response to or discussion of other work \\
\bottomrule
\end{tabular}
\end{table}

\subsection{Type Definitions}

\begin{definition}[Empirical Study (EMP)]
\label{def:emp}
A publication presenting \textbf{new primary data} collected through observation, experiment, survey, or other empirical methods, with original analysis of that data.

\textbf{Characteristic claims}: ``We find that...'', ``Our results show...'', ``The data indicate...''

\textbf{Examples}: RCTs, laboratory experiments, field studies, survey research, observational studies.
\end{definition}

\begin{definition}[Meta-Analysis (META)]
\label{def:meta}
A publication providing \textbf{quantitative synthesis} of multiple primary studies, using statistical methods to aggregate effect sizes and assess heterogeneity.

\textbf{Characteristic claims}: ``Across N studies, the pooled effect is...'', ``Meta-analytic evidence shows...''

\textbf{Examples}: Cochrane reviews, IPD meta-analyses, network meta-analyses.
\end{definition}

\begin{definition}[Review / Survey (REV)]
\label{def:rev}
A publication providing \textbf{qualitative synthesis} of existing research, organizing and interpreting the literature without formal statistical aggregation.

\textbf{Characteristic claims}: ``The literature suggests...'', ``Research in this area has found...'', ``Several studies indicate...''

\textbf{Examples}: Narrative reviews, scoping reviews, systematic reviews without meta-analysis.
\end{definition}

\begin{definition}[Theoretical Paper (THE)]
\label{def:the}
A publication presenting a \textbf{new conceptual framework, formal model, or logical argument} without primary empirical data.

\textbf{Characteristic claims}: ``We propose that...'', ``It follows that...'', ``The model predicts...''

\textbf{Examples}: Game-theoretic models, axiomatic frameworks, philosophical arguments, mathematical proofs.
\end{definition}

\begin{definition}[Methods Paper (METH)]
\label{def:meth}
A publication introducing a \textbf{new methodology, statistical technique, measurement instrument, or analytical tool}.

\textbf{Characteristic claims}: ``The method achieves...'', ``Our approach outperforms...'', ``The instrument measures...''

\textbf{Examples}: New statistical tests, software tools, measurement scales, experimental protocols.
\end{definition}

\begin{definition}[Replication Study (REP)]
\label{def:rep}
A publication \textbf{reproducing the methods of a prior study} to assess the robustness of original findings.

\textbf{Characteristic claims}: ``We replicate/fail to replicate...'', ``The original finding is/is not robust...''

\textbf{Examples}: Direct replications, conceptual replications, registered replication reports.
\end{definition}

\begin{definition}[Simulation Study (SIM)]
\label{def:sim}
A publication using \textbf{computational models} to explore theoretical predictions, test assumptions, or generate synthetic data.

\textbf{Characteristic claims}: ``Simulations show...'', ``The model generates...'', ``Under conditions X, the system...''

\textbf{Examples}: Agent-based models, Monte Carlo simulations, system dynamics models.
\end{definition}

\begin{definition}[Case Study (CASE)]
\label{def:case}
A publication providing \textbf{in-depth analysis of one or few cases}, prioritizing depth over generalizability.

\textbf{Characteristic claims}: ``In this case...'', ``The patient presented...'', ``The organization demonstrated...''

\textbf{Examples}: Clinical case reports, organizational case studies, historical analyses.
\end{definition}

\begin{definition}[Commentary / Letter (COM)]
\label{def:com}
A publication providing \textbf{response, critique, or discussion} of other published work, without substantial new data or theory.

\textbf{Characteristic claims}: ``The authors overlook...'', ``An alternative interpretation...'', ``We note that...''

\textbf{Examples}: Letters to the editor, commentaries, corrections, replies.
\end{definition}

% =============================================================================
% F.3 REPLICATION APPLICABILITY
% =============================================================================

\section{Replication Probability Applicability}
\label{sec:replication}

A central question for ESL is: \textit{Which paper types have meaningful ``replication probability''?}

\begin{table}[h]
\centering
\caption{Replication Probability Applicability by Paper Type}
\label{tab:replication}
\begin{tabular}{lcl}
\toprule
\textbf{Type} & \textbf{R-Score?} & \textbf{Rationale} \\
\midrule
\textbf{EMP} & \checkmark Yes & Core case: empirical claims can be replicated \\
\textbf{META} & \checkmark Yes & Meta-analytic conclusions can be re-assessed \\
\textbf{REP} & \checkmark Yes & Replication studies are themselves replicable \\
\textbf{SIM} & $\sim$ Partial & Reproducible (same code), but ``replication'' differs \\
\textbf{CASE} & $\sim$ Partial & Similar cases can be examined, not strict replication \\
\textbf{REV} & $\times$ No & No original empirical claims to replicate \\
\textbf{THE} & $\times$ No & Validity via logic, not replication \\
\textbf{METH} & $\times$ No & Validated, not replicated \\
\textbf{COM} & $\times$ No & No standalone empirical claims \\
\bottomrule
\end{tabular}
\end{table}

\begin{keyinsight}[title={R-Score vs. K-Score}]
\textbf{R-Score} (Replication probability) applies only to types with empirical claims.

\textbf{K-Score} (Claim-evidence alignment) applies to \textbf{all} types---the evidence scale simply differs.

For non-empirical types, K measures calibration against type-appropriate evidence standards.
\end{keyinsight}

\subsection{Simulation Studies: A Special Case}

Simulation studies occupy an intermediate position:

\begin{itemize}
    \item \textbf{Reproducibility}: Given same code + parameters, results should be identical (deterministic) or statistically equivalent (stochastic). This is \textit{computational reproducibility}, not replication.

    \item \textbf{Robustness}: Results should hold under reasonable parameter variations. This is closer to conceptual replication.

    \item \textbf{Empirical grounding}: Claims about real-world phenomena based on simulations require empirical validation.
\end{itemize}

\begin{warningbox}[title={Simulation Claims Require Careful Classification}]
A simulation paper may contain:
\begin{enumerate}
    \item \textbf{Model claims} (``The model captures mechanism X'') $\rightarrow$ THE-type evaluation
    \item \textbf{Computational claims} (``The algorithm converges in N steps'') $\rightarrow$ METH-type evaluation
    \item \textbf{Empirical claims} (``The simulation predicts real-world outcome Y'') $\rightarrow$ EMP-type evaluation
\end{enumerate}
Each claim type requires its own E-scale.
\end{warningbox}

% =============================================================================
% F.4 TYPE-SPECIFIC EVIDENCE SCALES
% =============================================================================

\section{Type-Specific Evidence Scales}
\label{sec:evidence-scales}

This section provides E-scale calibrations for each paper type.

\subsection{Empirical Studies (EMP)}

Standard ESL evidence hierarchy applies (see Chapter 3). Summary:

\begin{table}[h]
\centering
\caption{EMP Evidence Scale (Standard ESL)}
\label{tab:emp-scale}
\begin{tabular}{lc}
\toprule
\textbf{Evidence Level} & \textbf{E-value} \\
\midrule
Systematic review with meta-analysis & 0.90--0.95 \\
Large-scale RCT / pre-registered experiment & 0.80--0.90 \\
Well-powered lab experiment & 0.65--0.75 \\
Observational study with controls & 0.50--0.65 \\
Cross-sectional survey & 0.40--0.55 \\
Pilot study / exploratory analysis & 0.25--0.40 \\
Anecdotal / author intuition & 0.05--0.20 \\
\bottomrule
\end{tabular}
\end{table}

\subsection{Meta-Analyses (META)}

Meta-analyses have their own evidence quality determinants:

\begin{table}[h]
\centering
\caption{META Evidence Scale}
\label{tab:meta-scale}
\begin{tabular}{lc}
\toprule
\textbf{Quality Indicator} & \textbf{E-value} \\
\midrule
Cochrane-standard with low heterogeneity ($I^2 < 25\%$) & 0.90--0.95 \\
Pre-registered, comprehensive search, low bias & 0.80--0.90 \\
Standard meta-analysis, moderate heterogeneity & 0.70--0.80 \\
High heterogeneity ($I^2 > 75\%$) or publication bias & 0.50--0.65 \\
Limited search, few studies ($k < 5$) & 0.35--0.50 \\
Vote-counting or informal aggregation & 0.20--0.35 \\
\bottomrule
\end{tabular}
\end{table}

\textbf{Key factors}: Search comprehensiveness, inclusion criteria clarity, heterogeneity ($I^2$), publication bias assessment (funnel plot, Egger's test), quality of included studies.

\subsection{Reviews (REV)}

Review quality depends on systematicity and source quality:

\begin{table}[h]
\centering
\caption{REV Evidence Scale}
\label{tab:rev-scale}
\begin{tabular}{lc}
\toprule
\textbf{Review Type / Quality} & \textbf{E-value} \\
\midrule
Systematic review with clear protocol & 0.75--0.85 \\
Comprehensive narrative review, multiple high-quality sources & 0.60--0.75 \\
Scoping review with defined methodology & 0.55--0.70 \\
Standard literature review, good coverage & 0.45--0.60 \\
Selective review, limited sources & 0.30--0.45 \\
Informal overview, author's perspective & 0.15--0.30 \\
\bottomrule
\end{tabular}
\end{table}

\textbf{Key factors}: Explicit search strategy, inclusion/exclusion criteria, source quality assessment, coverage of relevant literature, balanced presentation.

\subsection{Theoretical Papers (THE)}

Theoretical work is evaluated on logical rigor and internal consistency:

\begin{table}[h]
\centering
\caption{THE Evidence Scale (Proof Rigor)}
\label{tab:the-scale}
\begin{tabular}{lc}
\toprule
\textbf{Logical Status} & \textbf{E-value} \\
\midrule
Formal proof with verified axioms & 0.95--1.00 \\
Complete proof, standard mathematical rigor & 0.85--0.95 \\
Rigorous argument, minor gaps fillable & 0.75--0.85 \\
Plausible argument, some steps informal & 0.55--0.70 \\
Conceptual framework, logic outlined not proven & 0.40--0.55 \\
Intuitive argument, key steps missing & 0.20--0.40 \\
Speculation, no formal structure & 0.05--0.20 \\
\bottomrule
\end{tabular}
\end{table}

\textbf{Key factors}: Axiom clarity, proof completeness, logical validity, assumption transparency, connection to established theory.

\begin{examplebox}[title={Theoretical Paper Example}]
\textbf{Claim}: ``Nash equilibrium exists in all finite games'' (Nash, 1950)

\textbf{Boldness}: $B = 0.95$ (universal claim with ``all'')

\textbf{Evidence}: Complete mathematical proof using Brouwer fixed-point theorem. $E = 0.98$

\textbf{K-Score}: $K = 0.95/0.98 = 0.97$ (Highly Stable)

The strong claim is justified by rigorous proof.
\end{examplebox}

\subsection{Methods Papers (METH)}

Methods are evaluated on validation quality:

\begin{table}[h]
\centering
\caption{METH Evidence Scale (Validation)}
\label{tab:meth-scale}
\begin{tabular}{lc}
\toprule
\textbf{Validation Level} & \textbf{E-value} \\
\midrule
Multi-site validation, diverse benchmarks & 0.85--0.95 \\
Independent validation by other researchers & 0.75--0.85 \\
Comprehensive internal validation, simulations & 0.65--0.75 \\
Limited validation, single benchmark & 0.45--0.60 \\
Proof-of-concept only & 0.30--0.45 \\
No validation, theoretical description only & 0.10--0.25 \\
\bottomrule
\end{tabular}
\end{table}

\textbf{Key factors}: Benchmark diversity, comparison to alternatives, edge case testing, user studies (if applicable), reproducibility of implementation.

\subsection{Replication Studies (REP)}

Replication studies are evaluated similarly to EMP, with additional considerations:

\begin{table}[h]
\centering
\caption{REP Evidence Scale}
\label{tab:rep-scale}
\begin{tabular}{lc}
\toprule
\textbf{Replication Quality} & \textbf{E-value} \\
\midrule
Pre-registered, high-powered, multi-site & 0.90--0.95 \\
Pre-registered, well-powered single site & 0.80--0.90 \\
Direct replication, adequate power & 0.70--0.80 \\
Conceptual replication, good design & 0.60--0.75 \\
Under-powered replication attempt & 0.40--0.55 \\
Informal replication, significant deviations & 0.25--0.40 \\
\bottomrule
\end{tabular}
\end{table}

\textbf{Key factors}: Pre-registration, statistical power, fidelity to original protocol, sample characteristics, outcome measure alignment.

\subsection{Simulation Studies (SIM)}

\begin{table}[h]
\centering
\caption{SIM Evidence Scale (Computational Validity)}
\label{tab:sim-scale}
\begin{tabular}{lc}
\toprule
\textbf{Simulation Quality} & \textbf{E-value} \\
\midrule
Validated against empirical data, robust to parameters & 0.80--0.90 \\
Comprehensive sensitivity analysis, code available & 0.70--0.80 \\
Standard robustness checks, reasonable assumptions & 0.55--0.70 \\
Limited sensitivity analysis & 0.40--0.55 \\
Single parameter set, no robustness testing & 0.25--0.40 \\
No validation, illustrative only & 0.10--0.25 \\
\bottomrule
\end{tabular}
\end{table}

\textbf{Key factors}: Parameter sensitivity, assumption transparency, code availability, comparison to analytical solutions (if available), empirical grounding.

\subsection{Case Studies (CASE)}

\begin{table}[h]
\centering
\caption{CASE Evidence Scale}
\label{tab:case-scale}
\begin{tabular}{lc}
\toprule
\textbf{Case Study Quality} & \textbf{E-value} \\
\midrule
Multiple cases, systematic comparison, triangulation & 0.65--0.75 \\
Single rich case, multiple data sources & 0.50--0.65 \\
Detailed case, limited triangulation & 0.40--0.55 \\
Descriptive case, single data source & 0.25--0.40 \\
Anecdotal case report & 0.10--0.25 \\
\bottomrule
\end{tabular}
\end{table}

\textbf{Key factors}: Case selection rationale, data triangulation, rival explanation consideration, thick description, analytic generalization attempt.

\textbf{Note}: Case study E-values are inherently lower due to limited generalizability. This is appropriate---claims should be hedged accordingly ($B$ should be low).

\subsection{Commentaries (COM)}

\begin{table}[h]
\centering
\caption{COM Evidence Scale}
\label{tab:com-scale}
\begin{tabular}{lc}
\toprule
\textbf{Commentary Quality} & \textbf{E-value} \\
\midrule
Points to verifiable errors with evidence & 0.80--0.90 \\
Substantive critique with citations & 0.60--0.75 \\
Reasoned alternative interpretation & 0.45--0.60 \\
Opinion with some support & 0.30--0.45 \\
Unsupported opinion & 0.10--0.25 \\
\bottomrule
\end{tabular}
\end{table}

\textbf{Key factors}: Specificity of critique, evidential support, constructiveness, connection to broader literature.

% =============================================================================
% F.5 WORKFLOW INTEGRATION
% =============================================================================

\section{ESL Workflow Integration}
\label{sec:workflow}

Paper Type Classification becomes Step 0 in the ESL assessment workflow.

\subsection{Updated ESL Workflow}

\begin{table}[h]
\centering
\caption{Complete ESL Workflow with PTC}
\label{tab:workflow}
\begin{tabular}{clp{7cm}}
\toprule
\textbf{Step} & \textbf{Name} & \textbf{Description} \\
\midrule
0 & \textbf{PTC} & Classify paper type (EMP/META/REV/THE/METH/REP/SIM/CASE/COM) \\
1 & \textbf{CTC} & Claim-Type Classification (T1--T5) \\
2 & \textbf{B-Assessment} & Linguistic marker analysis for Boldness \\
3 & \textbf{E-Assessment} & Type-specific evidence evaluation \\
4 & \textbf{K-Calculation} & $K = \min(B,E)/\max(B,E)$ \\
5 & \textbf{Safeguards} & Apply PP, WLA, EVA, RAS \\
6 & \textbf{Aggregation} & Document-level K-score \\
\bottomrule
\end{tabular}
\end{table}

\subsection{PTC Decision Tree}

\begin{verbatim}
START
  │
  ├─ Does the paper present new primary data?
  │   ├─ YES: Is it reproducing a prior study?
  │   │        ├─ YES → REP (Replication)
  │   │        └─ NO → EMP (Empirical)
  │   │
  │   └─ NO: Does it aggregate existing studies quantitatively?
  │           ├─ YES → META (Meta-Analysis)
  │           └─ NO: Does it synthesize existing literature?
  │                   ├─ YES → REV (Review)
  │                   └─ NO: Is it primarily computational/simulation?
  │                           ├─ YES → SIM (Simulation)
  │                           └─ NO: Does it introduce new methodology?
  │                                   ├─ YES → METH (Methods)
  │                                   └─ NO: Is it formal theory/proof?
  │                                           ├─ YES → THE (Theoretical)
  │                                           └─ NO: Is it single-case analysis?
  │                                                   ├─ YES → CASE (Case Study)
  │                                                   └─ NO → COM (Commentary)
END
\end{verbatim}

\subsection{Mixed-Type Papers}

Many papers contain multiple types of content. For mixed papers:

\begin{principle}[Claim-Level Type Assignment]
\label{princ:claim-level}
When a paper contains multiple content types, assign paper type \textbf{at the claim level}, not the paper level. Each claim is evaluated against the E-scale appropriate to its type.
\end{principle}

\begin{examplebox}[title={Mixed-Type Paper Example}]
A paper might contain:
\begin{enumerate}
    \item ``We propose a new theoretical model...'' $\rightarrow$ THE, E via proof rigor
    \item ``We test the model with a laboratory experiment...'' $\rightarrow$ EMP, E via study design
    \item ``Simulations confirm the model's predictions...'' $\rightarrow$ SIM, E via computational validity
\end{enumerate}

Each claim gets its own K-score; document K is aggregated per Chapter 2 rules.
\end{examplebox}

% =============================================================================
% F.6 IMPLEMENTATION NOTES
% =============================================================================

\section{Implementation Notes}
\label{sec:implementation}

\subsection{LLM Prompt Integration}

For automated ESL assessment, include PTC as the first step:

\begin{verbatim}
STEP 0 - PAPER TYPE CLASSIFICATION:
Before analyzing claims, classify this paper's type:
- EMP: New primary data collection and analysis
- META: Quantitative synthesis of multiple studies
- REV: Qualitative literature synthesis
- THE: Theoretical framework or formal proof
- METH: New methodology or tool
- REP: Replication of prior study
- SIM: Computational modeling
- CASE: In-depth single-case analysis
- COM: Commentary or response

Primary type: [TYPE]
Secondary types (if mixed): [TYPES]

Select appropriate E-scale based on paper type.
\end{verbatim}

\subsection{Reporting Format}

ESL reports should include paper type:

\begin{verbatim}
═══════════════════════════════════════════════════════
ESL ASSESSMENT REPORT
═══════════════════════════════════════════════════════
Paper: [Title]
Type: EMP (Empirical Study)
───────────────────────────────────────────────────────
Claim 1: "Our results demonstrate that X causes Y"
  - Type: T1 (Empirical, EMP-assessed)
  - B = 0.88 (strong: "demonstrate")
  - E = 0.72 (lab experiment, n=200)
  - K = 0.82 [STABLE]
...
\end{verbatim}

\subsection{Type-Specific Thresholds}

Consider adjusting K-thresholds by paper type:

\begin{table}[h]
\centering
\caption{Suggested Type-Specific K-Thresholds}
\label{tab:thresholds}
\begin{tabular}{lccc}
\toprule
\textbf{Type} & \textbf{Stable} & \textbf{Partially Stable} & \textbf{Unstable} \\
\midrule
THE (Theoretical) & $\geq 0.85$ & $0.65$--$0.85$ & $< 0.65$ \\
EMP (Empirical) & $\geq 0.75$ & $0.50$--$0.75$ & $< 0.50$ \\
CASE (Case Study) & $\geq 0.70$ & $0.45$--$0.70$ & $< 0.45$ \\
COM (Commentary) & $\geq 0.65$ & $0.40$--$0.65$ & $< 0.40$ \\
\bottomrule
\end{tabular}
\end{table}

\textbf{Rationale}: Theoretical papers making strong logical claims should be held to higher standards (proofs should be rigorous). Case studies and commentaries inherently have lower E-ceilings, so thresholds are adjusted downward.

% =============================================================================
% F.7 SUMMARY
% =============================================================================

\section{Summary}
\label{sec:summary}

This appendix has extended ESL to accommodate the full range of scientific publication types:

\begin{enumerate}
    \item \textbf{Nine Paper Types}: EMP, META, REV, THE, METH, REP, SIM, CASE, COM---each with distinct claim structures and evidence standards.

    \item \textbf{R-Score Applicability}: Only EMP, META, and REP have traditional replication probability. Other types use K-scores with type-specific E-scales.

    \item \textbf{Type-Specific E-Scales}: Each paper type has calibrated evidence hierarchies reflecting domain-appropriate quality indicators.

    \item \textbf{Workflow Integration}: PTC becomes Step 0 in ESL assessment, determining which E-scale to apply.

    \item \textbf{Mixed-Type Handling}: Claims are typed individually, enabling appropriate evaluation within multi-method papers.
\end{enumerate}

\begin{eslbox}
\textbf{Appendix F Self-Assessment}

This appendix proposes a classification system ($B \approx 0.75$) based on established publication typologies and ESL principles ($E \approx 0.70$ for framework extension with face validity but limited empirical testing).

\textbf{Calculation}: $K = 0.70/0.75 = 0.93$

\textbf{Status}: \textcolor{eslgreen}{\textbf{Highly Stable}} ($K \geq 0.90$)

\textbf{Claim Types}:
\begin{itemize}
    \item Paper type taxonomy: T2 (Based on established typologies)
    \item E-scale values: T4 (Expert proposal, requires validation)
    \item Workflow integration: T3 (Logical extension of ESL)
\end{itemize}
\end{eslbox}

% =============================================================================
% REFERENCES
% =============================================================================

\newpage
\bibliographystyle{apalike}

\begin{thebibliography}{99}

\bibitem[Bem, 2004]{Bem2004}
Bem, D. J. (2004).
Writing the empirical journal article.
In J. M. Darley, M. P. Zanna, \& H. L. Roediger III (Eds.), \textit{The Compleat Academic} (2nd ed., pp. 185--219). APA.

\bibitem[Borenstein et al., 2009]{Borenstein2009}
Borenstein, M., Hedges, L. V., Higgins, J. P. T., \& Rothstein, H. R. (2009).
\textit{Introduction to Meta-Analysis}.
Wiley.

\bibitem[Grant \& Booth, 2009]{Grant2009}
Grant, M. J., \& Booth, A. (2009).
A typology of reviews: An analysis of 14 review types and associated methodologies.
\textit{Health Information \& Libraries Journal}, 26(2), 91--108.

\bibitem[Yin, 2018]{Yin2018}
Yin, R. K. (2018).
\textit{Case Study Research and Applications: Design and Methods} (6th ed.).
SAGE.

\end{thebibliography}

\end{document}
